\documentclass[12pt,a4paper]{article}

% -------------------------------------------------------
% Pakete
% -------------------------------------------------------
\usepackage[T1]{fontenc}
\usepackage[utf8]{inputenc}
\usepackage[ngerman]{babel}
%\usepackage{lmodern}
\usepackage[expansion=false,protrusion=true]{microtype}
\usepackage{geometry}
\geometry{left=3cm, right=2.5cm, top=3cm, bottom=3cm}

\usepackage{amsmath, amssymb, amsthm}
\usepackage{mathtools}
\usepackage{array}
\usepackage{booktabs}
\usepackage{longtable}
\usepackage{enumitem}
\usepackage{graphicx}
\usepackage{caption}
\usepackage{subcaption}
\usepackage{float}

% TikZ fuer ER-Diagramme
\usepackage{tikz}
\usetikzlibrary{shapes.geometric, arrows.meta, positioning, calc, decorations.pathmorphing}

% Listings fuer SQL
\usepackage{listings}
\usepackage{xcolor}

\lstdefinestyle{sqlstyle}{
  language=SQL,
  basicstyle=\ttfamily\small,
  keywordstyle=\color{blue!70!black}\bfseries,
  commentstyle=\color{gray}\itshape,
  stringstyle=\color{green!50!black},
  numbers=left,
  numberstyle=\tiny\color{gray},
  stepnumber=1,
  frame=single,
  rulecolor=\color{black!30},
  backgroundcolor=\color{gray!5},
  breaklines=true,
  showstringspaces=false,
  tabsize=2,
  captionpos=b,
  extendedchars=true
}
\lstset{style=sqlstyle}

% Hyperref
\usepackage[hidelinks, colorlinks=false]{hyperref}
\usepackage{url}

% Bibliographie
\usepackage[numbers,sort&compress]{natbib}

% Kopf-/Fusszeilen
\usepackage{fancyhdr}
\pagestyle{fancy}
\fancyhf{}
\fancyhead[L]{\small Tiefenorientierter Datenbankentwurf}
\fancyhead[R]{\small \thepage}
\fancyfoot[C]{\small \textit{Wissenschaftliche Arbeit -- Datenbankentwurf}}
\renewcommand{\headrulewidth}{0.4pt}
\renewcommand{\footrulewidth}{0.2pt}

% Theoreme
\theoremstyle{definition}
\newtheorem{definition}{Definition}[section]
\newtheorem{theorem}{Satz}[section]
\newtheorem{lemma}[theorem]{Lemma}
\newtheorem{corollary}[theorem]{Korollar}
\newtheorem{proposition}[theorem]{Proposition}
\theoremstyle{remark}
\newtheorem{remark}{Bemerkung}[section]
\newtheorem{example}{Beispiel}[section]

\newcommand{\R}{\mathcal{R}}
\newcommand{\att}{\mathit{attr}}
\newcommand{\FD}{\mathit{FD}}
\newcommand{\MVD}{\mathit{MVD}}
\newcommand{\depth}{\mathit{depth}}
\newcommand{\breadth}{\mathit{breadth}}
\newcommand{\cost}{\mathit{cost}}
\newcommand{\NF}{\mathit{NF}}

% -------------------------------------------------------
% Dokument
% -------------------------------------------------------
\begin{document}

\begin{titlepage}
  \centering
  \vspace*{2cm}
  {\LARGE\bfseries Tiefenorientierter Datenbankentwurf}
  \vspace{1.5cm}
  {\large Wissenschaftliche Arbeit zum relationalen Datenbankentwurf}
  \vspace{1.5cm}

  \rule{\textwidth}{0.4pt}\\[0.5em]
  \begin{tabular}{rl}
    \textbf{Themengebiet:} & Datenbanktheorie, Schemaentwurf, Normalisierung \\
    \textbf{Schwerpunkt:}  & Tiefenpfade vs.\ sternförmige Relationsstruktur \\
    \textbf{Autor:}        & Stephan Epp \\
    \textbf{Datum:}        & \today \\
  \end{tabular}
  \rule{\textwidth}{0.4pt}

  \vfill
  \textit{Abstract}\\[0.5em]
  \begin{minipage}{0.85\textwidth}
  \small
  Der vorliegende Beitrag untersucht ein fundamentales, jedoch in der Literatur
  bislang nur implizit behandeltes Prinzip des relationalen Datenbankentwurfs:
  die strukturelle Orientierung von Fremdschlüssel"|relationen in die
  \emph{Tiefe} (lineare oder baumartige Ketten) im Gegensatz zur
  \emph{Breite} (sternförmige Hub-Anordnungen). Es wird formal nachgewiesen,
  dass tiefe Relationenpfade in Bezug auf Normalisierungsgrad, funktionale
  Abhängigkeitsstruktur, Anomalievermeidung, Abfragekomplexität und
  Wartbarkeit gegenüber flachen Sternstrukturen überlegen sind. Die Argumentation
  stützt sich auf die klassische relationale Algebra nach Codd~\cite{Codd1970},
  die Normalformtheorie bis BCNF und 4NF~\cite{Codd1974, Fagin1977}, sowie
  jüngere Erkenntnisse aus der Datenbankoptimierung und dem Schema-Refactoring~\cite{Ambler2003}.
  \end{minipage}
\end{titlepage}

\tableofcontents
\newpage

% ==============================================================
\section{Einleitung}
\label{sec:einleitung}
% ==============================================================

Der Entwurf relationaler Datenbankschemata ist eine der zentralen Aufgaben der
Softwareentwicklung und der angewandten Informatik. Obwohl die theoretischen
Grundlagen -- insbesondere Codd's relationales Modell \cite{Codd1970} sowie die
daraus abgeleiteten Normalformen \cite{Codd1974, Date2003} -- seit Jahrzehnten
bekannt sind, zeigen praktische Systeme immer wieder charakteristische
Entwurfsfehler. Einer der häufigsten und gleichzeitig folgenreichsten Fehler ist
die sogenannte \emph{sternförmige Schemaanordnung}: Eine zentrale
,,Hub\grqq{}-Relation referenziert direkt eine Vielzahl unabhängiger
,,Spoke\grqq{}-Relationen über Fremdschlüssel, ohne dass die natürliche semantische
Tiefe der Domäne abgebildet wird.

Dieser Beitrag stellt die These auf und belegt sie formal:
\begin{quote}
\emph{Ein tiefenorientierter Datenbankentwurf -- charakterisiert durch lineare
oder baumartige Fremdschlüsselketten -- ist dem sternförmigen Entwurf in allen
wesentlichen Gütekriterien der Datenbanktheorie überlegen.}
\end{quote}

Die Argumentation gliedert sich wie folgt: Abschnitt~\ref{sec:grundlagen}
etabliert die formalen Grundlagen. Abschnitt~\ref{sec:topologien} definiert
die Topologien ,,Tiefe\grqq{} und ,,Breite\grqq{} präzise.
Abschnitt~\ref{sec:normalformen} zeigt den Zusammenhang mit der
Normalisierungstheorie. Abschnitt~\ref{sec:anomalien} behandelt
Anomalievermeidung, Abschnitt~\ref{sec:abfragen} die Abfragekomplexität,
Abschnitt~\ref{sec:wartbarkeit} Wartbarkeit und Erweiterbarkeit.
Abschnitt~\ref{sec:beweis} liefert den konsolidierten formalen Nachweis.
Abschnitt~\ref{sec:beispiel} illustriert die Theorie an einem durchgängigen
Praxisbeispiel. Abschnitt~\ref{sec:fazit} schließt mit einem Fazit.

% ==============================================================
\section{Formale Grundlagen}
\label{sec:grundlagen}
% ==============================================================

\subsection{Relationales Modell}

\begin{definition}[Relation]
Eine \emph{Relation} $R$ über einem Attributschema $\att(R) = \{A_1, A_2, \ldots, A_n\}$
ist eine endliche Menge von Tupeln $t: \att(R) \to \bigcup_{i} \mathrm{dom}(A_i)$,
wobei $t(A_i) \in \mathrm{dom}(A_i)$ für alle $i$ \cite{Codd1970}.
\end{definition}

\begin{definition}[Datenbankschema]
Ein \emph{Datenbankschema} $\mathcal{S} = (\mathcal{R}, \mathcal{F}, \mathcal{K})$
besteht aus einer endlichen Menge von Relationenschemata $\mathcal{R} = \{R_1, \ldots, R_m\}$,
einer Menge funktionaler Abhängigkeiten $\mathcal{F}$ sowie einer Menge von
Integritätsbedingungen $\mathcal{K}$ (insbesondere Primär- und Fremdschlüsselbedingungen)
\cite{Ullman1988}.
\end{definition}

\begin{definition}[Funktionale Abhängigkeit]
Sei $R$ ein Relationenschema und $X, Y \subseteq \att(R)$.
Die \emph{funktionale Abhängigkeit} $X \to Y$ gilt in $R$, wenn für alle
Instanzen $r$ von $R$ und alle Tupel $t_1, t_2 \in r$ gilt:
$t_1(X) = t_2(X) \implies t_1(Y) = t_2(Y)$ \cite{Armstrong1974}.
\end{definition}

\begin{definition}[Fremdschlüssel]
Sei $R_i, R_j \in \mathcal{R}$ mit $K_j \subseteq \att(R_j)$ dem Primärschlüssel
von $R_j$. Eine Attributmenge $F \subseteq \att(R_i)$ heißt
\emph{Fremdschlüssel in $R_i$ referenzierend $R_j$}, wenn für alle Datenbankinstanzen
$d$ gilt: $\pi_F(r_i) \subseteq \pi_{K_j}(r_j)$ \cite{Date2003}.
\end{definition}

\subsection{Referenzgraph}

\begin{definition}[Referenzgraph]
\label{def:refgraph}
Sei $\mathcal{S}$ ein Datenbankschema. Der \emph{Referenzgraph}
$G(\mathcal{S}) = (V, E)$ hat als Knotenmenge $V = \mathcal{R}$ und als
Kantenmenge $E = \{(R_i, R_j) \mid \exists\, F \subseteq \att(R_i)\colon
F \text{ ist FK in } R_i \text{ referenzierend } R_j\}$.
\end{definition}

Der Referenzgraph ist das zentrale Werkzeug zur Unterscheidung von
Tiefenstruktur und Breitenstruktur eines Schemas.

\subsection{Komplexitätsmaße}

\begin{definition}[Schemakomplexitätsmaße]
\label{def:masse}
Sei $G(\mathcal{S}) = (V, E)$ der Referenzgraph eines Schemas $\mathcal{S}$.
\begin{enumerate}[label=(\roman*)]
  \item Der \emph{maximale Grad} $\Delta(G) = \max_{v \in V} \deg(v)$ misst
        die maximale Anzahl von Fremdschlüsselreferenzen eines Knotens.
  \item Die \emph{maximale Pfadlänge} $\lambda(G) = \max_{u,v \in V} d(u,v)$
        (Diameter) misst die längste Referenzkette im Schema.
  \item Das \emph{Tiefe-Breite-Verhältnis} $\tau(G) = \lambda(G) / \Delta(G)$
        charakterisiert die strukturelle Orientierung des Schemas.
\end{enumerate}
\end{definition}

Ein Schema heißt \emph{tiefenorientiert}, wenn $\tau(G) \gg 1$, und
\emph{breitenorientiert (sternförmig)}, wenn $\tau(G) \ll 1$.

% ==============================================================
\section{Topologien: Tiefe vs.\ Breite}
\label{sec:topologien}
% ==============================================================

\subsection{Formale Charakterisierung}

\begin{definition}[Sternschema]
\label{def:stern}
Ein Schema $\mathcal{S}$ heißt \emph{$k$-Sternschema}, wenn es eine
,,Hub\grqq{}-Relation $H \in \mathcal{R}$ gibt mit $\deg(H) = k$ und
$\lambda(G(\mathcal{S})) \leq 2$. Die $k$ von $H$ referenzierten Relationen
heißen \emph{Spoke}-Relationen.
\end{definition}

Das Sternschema ist insbesondere aus dem Data-Warehouse-Bereich bekannt
(Kimball-Methodik, \cite{Kimball1996}), wo es für OLAP-Abfragen optimiert
wird. Im OLTP-Kontext hingegen offenbart es gravierende Schwächen
(vgl.\ Abschnitt~\ref{sec:anomalien}).

\begin{definition}[Tiefenpfadschema]
\label{def:tiefenpfad}
Ein Schema $\mathcal{S}$ heißt \emph{Tiefenpfadschema der Länge $l$},
wenn $G(\mathcal{S})$ einen Pfad $R_1 \to R_2 \to \cdots \to R_l$
enthält mit $\Delta(G) \leq c$ für eine kleine Konstante $c$ (typisch $c \leq 3$)
und $\lambda(G(\mathcal{S})) = l - 1$.
\end{definition}

\begin{figure}[H]
\centering
\begin{tikzpicture}[
  node distance=1.4cm and 2.0cm,
  box/.style={rectangle, draw, rounded corners=3pt, minimum width=2.2cm,
              minimum height=0.8cm, font=\small\ttfamily, fill=blue!8},
  hub/.style={rectangle, draw, rounded corners=3pt, minimum width=2.5cm,
              minimum height=0.9cm, font=\small\ttfamily\bfseries, fill=orange!20},
  arr/.style={-{Stealth[length=7pt]}, thick},
  darr/.style={-{Stealth[length=7pt]}, thick, blue!70}
]

% === Sternschema (links) ===
\node[hub] (H) at (0,0) {Hub $H$};
\foreach \i/\name/\angle in {
  1/Spoke$_1$/90,
  2/Spoke$_2$/30,
  3/Spoke$_3$/-30,
  4/Spoke$_4$/-90,
  5/Spoke$_5$/150,
  6/Spoke$_6$/210
}{
  \node[box] (S\i) at ([shift={(\angle:2.8cm)}]H) {\name};
  \draw[arr, orange!60!black] (H) -- (S\i);
}
\node[above=0.15cm of H, font=\small\bfseries] {(a) Sternschema};

% === Tiefenpfadschema (rechts) ===
\begin{scope}[xshift=8.5cm]
  \node[box] (R1) at (0, 2.4)   {$R_1$};
  \node[box] (R2) at (0, 1.0)   {$R_2$};
  \node[box] (R3) at (0,-0.4)   {$R_3$};
  \node[box] (R4) at (0,-1.8)   {$R_4$};
  \node[box] (R5) at (0,-3.2)   {$R_5$};

  \draw[darr] (R1) -- (R2);
  \draw[darr] (R2) -- (R3);
  \draw[darr] (R3) -- (R4);
  \draw[darr] (R4) -- (R5);

  \node[above=0.15cm of R1, font=\small\bfseries] {(b) Tiefenpfadschema};
\end{scope}

\end{tikzpicture}
\caption{Gegenüberstellung von (a) Sternschema mit Hub-Relation $H$
und 6 Spoke-Relationen ($\Delta = 6$, $\lambda = 1$, $\tau = 0{,}17$) und
(b) Tiefenpfadschema mit linearer Kette
($\Delta = 1$, $\lambda = 4$, $\tau = 4$).}
\label{fig:topologien}
\end{figure}

\subsection{Semantische Motivation}

Die Wahl der Topologie ist kein rein ästhetisches, sondern ein
\emph{semantisch erzwungenes} Entwurfsprinzip. Reale Domänen weisen
fast immer hierarchische oder transitive Abhängigkeitsstrukturen auf:
Eine \texttt{Bestellung} hängt von einem \texttt{Kunden} ab, der wiederum
einer \texttt{Region} zugeordnet ist, die einem \texttt{Land} angehört.
Diese natürliche Tiefe der Abhängigkeiten erzwingt eine tiefe Pfadstruktur
und sollte nicht künstlich in eine flache Sternstruktur gezwungen werden
\cite{Date2003, Elmasri2015}.

% ==============================================================
\section{Normalisierungstheorie und Schematopologie}
\label{sec:normalformen}
% ==============================================================

\subsection{Normalformen und ihre Anforderungen}

Wir rekapitulieren die für den Nachweis relevanten Normalformen \cite{Codd1974, Fagin1977}:

\begin{definition}[Erste Normalform -- 1NF]
Eine Relation $R$ befindet sich in \emph{1NF}, wenn alle Attributdomänen
atomar sind, d.\,h.\ keine mengenwertigen oder strukturierten Attribute
enthält \cite{Codd1970}.
\end{definition}

\begin{definition}[Zweite Normalform -- 2NF]
Eine Relation $R$ in 1NF befindet sich in \emph{2NF}, wenn jedes
Nichtschlüsselattribut voll funktional abhängig vom Primärschlüssel ist,
d.\,h.\ keine partielle Abhängigkeit von einem echten Teilschlüssel existiert
\cite{Codd1974}.
\end{definition}

\begin{definition}[Dritte Normalform -- 3NF]
Eine Relation $R$ in 2NF befindet sich in \emph{3NF}, wenn kein
Nichtschlüsselattribut transitiv vom Primärschlüssel abhängig ist
\cite{Codd1974}.
\end{definition}

\begin{definition}[Boyce-Codd-Normalform -- BCNF]
Eine Relation $R$ befindet sich in \emph{BCNF}, wenn für jede
nicht-triviale funktionale Abhängigkeit $X \to A$ gilt: $X$ ist ein
Superschlüssel von $R$ \cite{Boyce1974}.
\end{definition}

\begin{definition}[Vierte Normalform -- 4NF]
Eine Relation $R$ in BCNF befindet sich in \emph{4NF}, wenn für jede
nicht-triviale mehrwertige Abhängigkeit $X \twoheadrightarrow Y$ gilt:
$X$ ist ein Superschlüssel von $R$ \cite{Fagin1977}.
\end{definition}

\subsection{Transitivität und Schematopologie}

\begin{theorem}[Transitivitätsverteilung]
\label{thm:transitivitaet}
Sei $R$ eine Relation mit $\att(R) = \{K, A_1, A_2, \ldots, A_n\}$ und
funktionalen Abhängigkeiten $K \to A_1 \to A_2 \to \cdots \to A_n$.
Dann verletzt $R$ die 3NF für alle $A_i$ mit $i \geq 2$.
Die einzige verlustfreie, abhängigkeitserhaltende Zerlegung in 3NF ist
die Kette $R_1(K, A_1),\; R_2(A_1, A_2),\; \ldots,\; R_{n-1}(A_{n-1}, A_n)$.
\end{theorem}

\begin{proof}
Sei $K \to A_i \to A_{i+1}$ für ein $i \geq 1$ eine transitive Kette.
Da $A_i$ kein Schlüssel ist, ist $A_{i+1}$ transitiv von $K$ abhängig und
$R$ verletzt 3NF. Die Zerlegung $R' = R \setminus \{A_{i+1}\}$ und
$R'' = (A_i, A_{i+1})$ ist verlustfrei nach dem Lossless-Join-Lemma
(da $A_i \to A_{i+1}$ gilt und $A_i = R' \cap R''$ sowie $A_i \to R''$ gilt)
\cite{Ullman1988}. Rekursive Anwendung für alle transitiven Abhängigkeiten
liefert die vollständige Zerlegungskette. Die Abhängigkeitserhaltung folgt
daraus, dass jede FD $A_i \to A_{i+1}$ in der entsprechenden Relation
$R_i(A_{i-1}, A_i, A_{i+1})$ bzw.\ $R_i(A_i, A_{i+1})$ präsent bleibt.
\end{proof}

\begin{corollary}[Tiefenpfad als 3NF-Normalform]
\label{cor:tiefenpfad3NF}
Ein Datenbankschema, das alle transitiven Abhängigkeiten einer Domäne
korrekt repräsentiert, hat zwingend einen tiefenorientierten Referenzgraphen.
Sternförmige Schemata, die dieselbe Domäne abbilden, verletzen die 3NF,
sofern transitive Abhängigkeiten in der Domäne existieren.
\end{corollary}

Korollar~\ref{cor:tiefenpfad3NF} ist das zentrale formale Ergebnis dieses
Abschnitts. Es besagt, dass die Normalisierungstheorie \emph{tiefe Pfade}
direkt erzwingt, wenn die Domäne transitive Abhängigkeiten aufweist --
was in praktisch allen nicht-trivialen realen Domänen der Fall ist.

\subsection{Mehrwertige Abhängigkeiten und 4NF}

\begin{proposition}[Mehrwertige Abhängigkeiten im Sternschema]
\label{prop:mvd_stern}
In einem $k$-Sternschema mit Hub $H$ und Spokes $S_1, \ldots, S_k$ gilt:
Falls zwei Spoke-Attribute $B_1 \in \att(S_1)$ und $B_2 \in \att(S_2)$
inhaltlich unabhängig sind und beide durch $K_H$ identifiziert werden, so
liegt in $H$ eine nicht-triviale mehrwertige Abhängigkeit $K_H \twoheadrightarrow B_1$
vor, wenn beide Attribute direkt in $H$ gehalten werden. Dies verletzt 4NF.
\end{proposition}

Die korrekte Auflösung ist -- erneut -- die Extraktion in separate Relationen,
was die Tiefenstruktur des Schemas weiter verstärkt.

% ==============================================================
\section{Anomalievermeidung}
\label{sec:anomalien}
% ==============================================================

\subsection{Klassifikation von Anomalien}

Die Datenbanktheorie unterscheidet drei klassische Anomalietypen \cite{Date2003}:

\begin{enumerate}[label=(\roman*)]
  \item \textbf{Einfügeanomalie}: Daten können nicht eingetragen werden,
        ohne gleichzeitig nicht vorhandene Daten erfinden zu müssen.
  \item \textbf{Löschanomalie}: Das Löschen eines Tupels vernichtet
        unbeabsichtigt weitere Informationen.
  \item \textbf{Änderungsanomalie}: Eine inhaltliche Änderung muss in
        mehreren Tupeln durchgeführt werden, da die Information redundant gespeichert ist.
\end{enumerate}

\begin{theorem}[Anomaliefreiheit durch Tiefenpfade]
\label{thm:anomaliefrei}
Sei $\mathcal{S}$ ein Datenbankschema in BCNF mit einem Tiefenpfadgraphen
$G(\mathcal{S})$. Dann ist $\mathcal{S}$ frei von Einfüge-, Lösch- und
Änderungsanomalien bzgl.\ der durch $G(\mathcal{S})$ modellierten
transitiven Abhängigkeiten.
\end{theorem}

\begin{proof}
BCNF garantiert, dass in jeder Relation $R_i$ der Kette jede nicht-triviale
FD $X \to A$ mit $X$ als Superschlüssel gilt. Somit ist jede Information
genau einmal gespeichert (keine Redundanz). Ohne Redundanz:
\begin{itemize}
  \item \emph{Änderungsanomalie}: Entfällt, da jede Eigenschaft in genau einer
        Relation steht. Eine Änderung betrifft genau ein Tupel.
  \item \emph{Löschanomalie}: Da eine Entität $e_i$ in $R_i$ unabhängig von
        ihrer Referenz durch $R_{i-1}$ existiert (referentielle Integrität
        via FK), gehen beim Löschen eines $R_{i-1}$-Tupels keine $R_i$-Daten
        verloren (ON DELETE SET NULL / RESTRICT).
  \item \emph{Einfügeanomalie}: Eine neue Entität in $R_i$ kann unabhängig
        von $R_{i-1}$ eingefügt werden, da FK-Referenz nur in der
        referenzierenden Relation besteht.
\end{itemize}
Im Sternschema hingegen erzeugt die direkte Kodierung transitiver Attribute
in $H$ Redundanz, da Attributwerte von Spoke-Entitäten wiederholt in $H$
gespeichert sind -- dies ermöglicht alle drei Anomalietypen.
\end{proof}

\subsection{Quantitative Analyse der Redundanz}

\begin{proposition}[Redundanzfaktor im Sternschema]
\label{prop:redundanz}
Sei $H$ eine Hub-Relation mit $n_H$ Tupeln und Spoke-Relation $S_1$ mit
$n_{S_1}$ Tupeln und $|A_{S_1}|$ Attributen. Werden die Attribute von $S_1$
direkt in $H$ denormalisiert, so beträgt der Redundanzfaktor:
\[
  \rho = \frac{n_H}{n_{S_1}} \cdot |A_{S_1}|
\]
Für $n_H \gg n_{S_1}$ (typisch bei 1:n-Beziehungen) gilt $\rho \gg 1$.
\end{proposition}

In einer typischen Bestelldatenbank mit $n_H = 10^6$ Bestellungen und
$n_{S_1} = 10^4$ Kunden ergibt sich $\rho = 100 \cdot |A_{\text{Kunde}}|$.
Bei 10 Kundenattributen werden also 1\,000-fach mehr Daten gespeichert
als notwendig -- ein massiver Speicher- und Konsistenzaufwand.

% ==============================================================
\section{Abfragekomplexität und Optimierbarkeit}
\label{sec:abfragen}
% ==============================================================

\subsection{Join-Komplexität}

Ein häufiges Argument gegen tiefe Pfade ist der höhere Join-Aufwand.
Wir zeigen, dass dieses Argument zwar formal korrekt, aber praktisch
irrelevant ist, wenn moderne Datenbankoptimierung berücksichtigt wird.

\begin{definition}[Join-Tiefe]
Die \emph{Join-Tiefe} $\delta(\mathcal{Q})$ einer Abfrage $\mathcal{Q}$ über
Schema $\mathcal{S}$ ist die Anzahl der JOIN-Operationen in der optimalen
logischen Ausführung von $\mathcal{Q}$.
\end{definition}

\begin{proposition}[Join-Tiefe und Pfadlänge]
\label{prop:jointiefe}
Für eine vollständige Traversal-Abfrage über einen Tiefenpfad
$R_1 \to R_2 \to \cdots \to R_l$ gilt $\delta = l - 1$.
Für ein $k$-Sternschema gilt für dieselbe semantische Abfrage $\delta = k$.
Da $l \ll k$ in tiefen Schemata (wenige, längere Ketten statt vieler
direkter Abhängigkeiten), ist die Join-Tiefe im Tiefenpfad geringer
für partielle Traversals \cite{Ramakrishnan2003}.
\end{proposition}

\begin{remark}
Selbst wenn $l > k$ für globale Traversals gilt, kompensiert dies der
Abfrageoptimierer durch:
\begin{itemize}
  \item \emph{Index-Only Scans} auf FK-Spalten (eliminiert Heap-Zugriffe),
  \item \emph{Hash-Joins} mit amortisierter $O(n)$-Komplexität \cite{Graefe1993},
  \item \emph{Materialized Views} für häufige tiefe Traversals,
  \item \emph{Predicate Pushdown} entlang des Pfades.
\end{itemize}
\end{remark}

\subsection{Selektivität und Filtereffizienz}

\begin{theorem}[Selektivitätsvorteil tiefer Pfade]
\label{thm:selektivitaet}
Sei $\mathcal{Q}$ eine Abfrage mit Filterprädikat $\sigma_P$ auf Attributen
einer Entität $E$ der Tiefe $d$ im Pfad. Im Tiefenpfadschema kann $\sigma_P$
direkt auf der Relation $R_d$ angewandt werden (Predicate Pushdown).
Im Sternschema hingegen sind transitive Attribute von $E$ in der Hub-Relation
$H$ de\-normalisiert, was bedeutet, dass der Filter auf der gesamten,
$n_H$-tupligen Hub-Relation angewandt werden muss.
\[
  \text{Kosten}_{\text{Stern}}(\sigma_P) = O(n_H) \cdot c_{\text{IO}}
  \qquad\text{vs.}\qquad
  \text{Kosten}_{\text{Tief}}(\sigma_P) = O(n_{R_d}) \cdot c_{\text{IO}}
\]
Da $n_H \geq n_{R_d}$ für alle $d$, ist der Tiefenpfad mindestens so
effizient wie das Sternschema, typisch deutlich effizienter.
\end{theorem}

\begin{proof}
Im normalisierten Tiefenpfadschema liegen die Attribute von Entität $E$ in
$R_d$ mit $|R_d| = n_{R_d}$ Tupeln. Ein Index auf dem Primärschlüssel von
$R_d$ ermöglicht $O(\log n_{R_d} + k')$ für $k'$ Ergebnistupel.
Im Sternschema befinden sich diese Attribute in $H$ vermischt mit allen anderen
Attributen; ohne Partial Index muss der gesamte Heap von $H$ gescannt werden
($O(n_H)$). Da eine 1:n-Beziehung $|E| \leq n_H$ impliziert, gilt
$n_{R_d} \leq n_H$, und die Ungleichung ist bewiesen.
\end{proof}

\subsection{Schreibkomplexität und Konsistenzerhaltung}

\begin{proposition}[Schreibkostenvorteil]
Im Tiefenpfadschema erfordert eine Änderung an Attributen der Entität $E$
genau \emph{eine} UPDATE-Operation auf $R_d$. Im denormalisierten
Sternschema müssen alle Tupel in $H$, die auf diese Entität verweisen,
aktualisiert werden: im Worst Case $O(n_H / n_{R_d})$ Operationen --
proportional zum Redundanzfaktor $\rho$ aus Proposition~\ref{prop:redundanz}.
\end{proposition}

% ==============================================================
\section{Wartbarkeit und Erweiterbarkeit}
\label{sec:wartbarkeit}
% ==============================================================

\subsection{Schemaevolution}

Ein zentrales Qualitätsmerkmal eines Datenbankschemas ist seine
Fähigkeit, sich an veränderte Anforderungen anzupassen -- bekannt als
\emph{Schemaevolution} \cite{Roddick1995}.

\begin{theorem}[Lokale Änderbarkeit im Tiefenpfad]
\label{thm:schaemaevolution}
Im Tiefenpfadschema $R_1 \to R_2 \to \cdots \to R_l$ ist das Hinzufügen,
Entfernen oder Ändern von Attributen der Entität $R_i$ eine
\emph{lokale Schemaoperation}, die ausschließlich $R_i$ betrifft.
Im $k$-Sternschema mit denormalisierten Attributen in $H$ ist dieselbe
Operation \emph{global} und betrifft $H$ sowie potenziell alle
abhängigen Abfragen, Views und Applikationsschichten.
\end{theorem}

\begin{proof}
Da im Tiefenpfadschema jedes Attribut in genau einer Relation vorkommt
(BCNF-Eigenschaft), hat eine Schemaänderung an $R_i$ keinen direkten
Einfluss auf $R_j$ mit $j \neq i$. Nur FK-Referenzen
$R_{i-1} \to R_i$ sind betroffen -- diese beschränken sich auf eine
einzige Kante im Referenzgraphen. Im Sternschema sind Attribute aus
multiple Entitäten in $H$ gemischt; eine Änderung an $H$ betrifft
alle darauf basierenden Abfragen, was zu einem kaskadierten
Änderungsaufwand führt \cite{Ambler2003}.
\end{proof}

\subsection{Zugriffsrechteverwaltung}

Ein weiterer Vorteil tiefer Pfade zeigt sich bei der Granularität der
Zugriffsrechteverwaltung:

\begin{proposition}[Berechtigungsgranularität]
Im Tiefenpfadschema können Zugriffsrechte \emph{entitätspräzise} vergeben
werden: Ein Nutzer erhält SELECT-Rechte auf $R_d$, ohne Zugriff auf andere
Relationen zu erhalten. Im Sternschema mit Hub $H$ bedeutet Zugriff auf
$H$ stets Zugriff auf alle in $H$ denormalisierten Entitätsdaten --
eine grobe Granularität, die dem Prinzip der minimalen Rechtevergabe
(Principle of Least Privilege) widerspricht \cite{Saltzer1975}.
\end{proposition}

\subsection{Testbarkeit und Datenqualität}

Tiefenpfadschemata sind inhärent leichter zu testen und zu validieren,
da:
\begin{itemize}
  \item Jede Relation eine \emph{klar umgrenzte semantische Einheit} darstellt
        (Single Responsibility für Relationen),
  \item Integritätsbedingungen pro Relation formuliert und geprüft werden,
  \item Fremdschlüsselbedingungen entlang des Pfades eine implizite
        Validierungskaskade bilden.
\end{itemize}

% ==============================================================
\section{Konsolidierter formaler Nachweis}
\label{sec:beweis}
% ==============================================================

\subsection{Gütekriterien}

Wir formalisieren nun den umfassenden Vergleich anhand eines Vektors von
Gütekriterien:

\begin{definition}[Gütevektor]
Sei $\mathcal{S}$ ein Datenbankschema. Der \emph{Gütevektor}
$\mathbf{q}(\mathcal{S}) \in \mathbb{R}^6$ ist definiert als:
\[
  \mathbf{q}(\mathcal{S}) = \bigl(
    q_{\NF},\; q_{\text{Red}},\; q_{\text{Anom}},\;
    q_{\text{Sch}},\; q_{\text{Wart}},\; q_{\text{Sek}}
  \bigr)
\]
wobei die Komponenten Normalformgrad, Redundanzfreiheit,
Anomaliefreiheit, Schreibkosteneffizienz, Wartbarkeit und
Selektivitätseffizienz kodieren (alle normiert auf $[0,1]$,
höher ist besser).
\end{definition}

\begin{theorem}[Ãœberlegenheit des Tiefenpfadschemas]
\label{thm:hauptsatz}
Sei $\mathcal{S}_{\text{Tief}}$ ein Tiefenpfadschema in BCNF und
$\mathcal{S}_{\text{Stern}}$ ein semantisch äquivalentes $k$-Sternschema
($k \geq 3$) mit denormalisierten transitiven Attributen. Dann gilt:
\[
  \mathbf{q}(\mathcal{S}_{\text{Tief}}) \succ \mathbf{q}(\mathcal{S}_{\text{Stern}})
\]
komponentenweise, d.\,h.\ $q_j(\mathcal{S}_{\text{Tief}}) \geq q_j(\mathcal{S}_{\text{Stern}})$
für alle $j$, mit strikter Ungleichung für mindestens vier Komponenten.
\end{theorem}

\begin{proof}
Wir zeigen die strikten Ungleichungen komponentenweise:

\medskip
(i) \textbf{Normalformgrad} $q_{\NF}$: Nach Korollar~\ref{cor:tiefenpfad3NF}
befindet sich $\mathcal{S}_{\text{Tief}}$ in BCNF, während
$\mathcal{S}_{\text{Stern}}$ transitive Abhängigkeiten in $H$ enthält
und damit 3NF verletzt. Somit $q_{\NF}(\mathcal{S}_{\text{Tief}}) > q_{\NF}(\mathcal{S}_{\text{Stern}})$.

(ii) \textbf{Redundanzfreiheit} $q_{\text{Red}}$: Nach Proposition~\ref{prop:redundanz}
ist $\rho(\mathcal{S}_{\text{Stern}}) \gg 1$, während
$\rho(\mathcal{S}_{\text{Tief}}) = 1$ (jede Information einmalig gespeichert).

(iii) \textbf{Anomaliefreiheit} $q_{\text{Anom}}$: Folgt direkt aus
Satz~\ref{thm:anomaliefrei}.

(iv) \textbf{Schreibkosteneffizienz} $q_{\text{Sch}}$: Aus der Redundanz
$\rho > 1$ im Sternschema folgt ein entsprechend höherer Schreibaufwand.

(v) \textbf{Wartbarkeit} $q_{\text{Wart}}$: Nach Satz~\ref{thm:schaemaevolution}
sind Schemaänderungen im Tiefenpfad lokal beschränkt.

(vi) \textbf{Selektivitätseffizienz} $q_{\text{Sek}}$: Nach
Satz~\ref{thm:selektivitaet} gilt stets
$\text{Kosten}_{\text{Tief}} \leq \text{Kosten}_{\text{Stern}}$.

Da für (i)--(v) strikte Ungleichungen gelten, ist die Behauptung bewiesen.
\end{proof}

% ==============================================================
\section{Graphentheoretische Fundierung: DFS, BFS und die Asymmetrie der Suchalgorithmen}
\label{sec:dfs_bfs}
% ==============================================================

Die bisherigen Argumente stützten sich auf die Normalformtheorie und die
Anomalielehre. Es gibt jedoch eine tiefere, algorithmische Begründung für
die Ãœberlegenheit tiefer Strukturen, die in der Theorie der Graphentraversierung
verwurzelt ist. Die Beobachtung ist fundamental:

\begin{quote}
\emph{Tiefensuche (DFS) und Breitensuche (BFS) über den Referenzgraphen eines
Datenbankschemas sind \textbf{nicht symmetrisch}. Ihre Ergebnisstrukturen --
und damit die induzierten Schemastrukturen -- unterscheiden sich qualitativ
und in ihrer Komplexität grundlegend.}
\end{quote}

\subsection{Tiefensuche und der DFS-Wald}

\begin{definition}[DFS-Baum und DFS-Wald, nach Cormen et al.~\cite{Cormen2009}]
\label{def:dfswald}
Sei $G = (V, E)$ ein gerichteter Graph. Die \emph{Tiefensuche} (Depth-First
Search, DFS) traversiert $G$, indem sie von jedem besuchten Knoten $u$ aus
rekursiv in die Nachfolger abtaucht, bevor sie zum Ausgangspunkt zurückkehrt.
Die dabei entstehenden Baumkanten bilden einen \emph{DFS-Wald}
$T_{\mathrm{DFS}} = (V, E_{\mathrm{tree}})$ mit
$E_{\mathrm{tree}} \subseteq E$. Jede Zusammenhangskomponente des DFS-Walds
ist ein \emph{DFS-Baum}.
\end{definition}

Der DFS-Wald hat bemerkenswerte Eigenschaften, die ihn für den
Datenbankschemaentwurf auszeichnen:

\begin{theorem}[Struktureigenschaften des DFS-Walds, \cite{Cormen2009, Tarjan1972}]
\label{thm:dfswald}
Sei $T_{\mathrm{DFS}}$ der DFS-Wald über dem Referenzgraphen
$G(\mathcal{S})$. Dann gilt:
\begin{enumerate}[label=(\roman*)]
  \item $T_{\mathrm{DFS}}$ ist ein \emph{Wald} (azyklischer Graph, Menge
        von Bäumen) -- also eine azyklische, hierarchische Struktur.
  \item Die Tiefe $\depth(T_{\mathrm{DFS}})$ ist maximal unter allen
        aufspannenden Wäldern von $G(\mathcal{S})$.
  \item Rückwärtskanten (back edges) $e \notin E_{\mathrm{tree}}$ identifizieren
        Zyklen und damit potenzielle \emph{zirkuläre Abhängigkeiten} im Schema.
  \item Die \emph{topologische Sortierung} des Schemas -- notwendig für
        eine anomaliefreie Einfüge- und Löschreihenfolge -- ist
        genau die DFS-Postorder-Reihenfolge.
\end{enumerate}
\end{theorem}

\begin{corollary}[DFS-Wald als natürliches Schemaabbild]
Der DFS-Wald $T_{\mathrm{DFS}}$ über $G(\mathcal{S})$ ist das
\emph{natürliche strukturelle Abbild} eines korrekt normalisierten
Tiefenpfadschemas. Die Baumkanten entsprechen Fremdschlüsselbeziehungen,
die Baumtiefe der semantischen Abhängigkeitskette.
Jedes BCNF-Schema erzeugt -- wenn als Graph betrachtet -- einen azyklischen,
tiefenorientierten Referenzgraphen, der einem DFS-Wald isomorph ist.
\end{corollary}

\subsection{Breitensuche und das Fehlen eines BFS-Walds}

Die Breitensuche (Breadth-First Search, BFS) traversiert den Graphen
schichtweise: Alle Knoten der Distanz $d$ werden vollständig besucht, bevor
Knoten der Distanz $d+1$ betreten werden.

\begin{definition}[BFS-Baum, nach Cormen et al.~\cite{Cormen2009}]
\label{def:bfsbaum}
Die Breitensuche über $G = (V, E)$ ausgehend von einem Startknoten $s$
erzeugt einen \emph{BFS-Baum} $T_{\mathrm{BFS}}$, der die kürzesten Pfade
von $s$ zu allen erreichbaren Knoten enthält. Die Breite des BFS-Baums auf
Tiefe $d$ ist $b_d = |\{v \mid d_G(s,v) = d\}|$.
\end{definition}

\begin{theorem}[BFS erzeugt keinen Wald]
\label{thm:bfskeinwald}
Die Breitensuche über einem unzusammenhängenden Graphen $G$ mit $k$
Zusammenhangskomponenten $C_1, \ldots, C_k$ erzeugt im Allgemeinen
\emph{keinen Wald}, sondern eine Menge von $k$ BFS-Bäumen
$T_{\mathrm{BFS}}^{(1)}, \ldots, T_{\mathrm{BFS}}^{(k)}$, die jeweils
\emph{sternförmig} um ihren Startknoten $s_i \in C_i$ angeordnet sind.
Insbesondere:
\begin{enumerate}[label=(\roman*)]
  \item Für einen $k$-Stern\-graph $G_\star$ mit Hub $H$ und Knoten
        $S_1, \ldots, S_k$ ist der BFS-Baum mit Start $H$ identisch
        mit $G_\star$ selbst: $T_{\mathrm{BFS}} = G_\star$.
  \item Die BFS-Tiefe beträgt $\depth(T_{\mathrm{BFS}}) = 1$, unabhängig
        von $k$.
  \item Die BFS über einem zusammenhängenden Graphen mit $n$ Knoten
        erzeugt genau \emph{einen} Baum, nicht einen Wald -- der Begriff
        ,,BFS-Wald\grqq{} ist daher für zusammenhängende Graphen
        nicht definiert.
\end{enumerate}
\end{theorem}

\begin{proof}
(i) Beim Startsknoten $H$ werden in der ersten BFS-Schicht alle direkten
Nachbarn $S_1, \ldots, S_k$ besucht. Da diese keine weiteren Nachbarn
besitzen (Stern-Definition), terminiert BFS nach Schicht 1.
Die Baumkanten sind genau $\{(H, S_i) \mid i = 1, \ldots, k\}$, was dem
Stern entspricht.

(ii) $\depth(T_{\mathrm{BFS}}) = \max_{v} d_{T}(s, v) = 1$.

(iii) BFS ist eine \emph{Single-Source}-Traversierung; bei zusammenhängendem
$G$ wird genau ein Baum mit Wurzel $s$ erzeugt. Mehrere Wurzeln (Wald)
entstehen nur durch Multi-Source-BFS, was der Standard-BFS-Definition
widerspricht. Im Gegensatz dazu erzeugt DFS über einem unzusammenhängenden
Graphen automatisch einen Wald, indem für jede noch unbesuchte Zusammenhangskomponente
eine neue DFS-Traversierung gestartet wird.
\end{proof}

\subsection{Die fundamentale Asymmetrie von DFS und BFS}

\begin{theorem}[Asymmetrie von DFS und BFS]
\label{thm:asymmetrie}
DFS und BFS sind \emph{algorithmisch asymmetrisch} in folgenden Dimensionen:
\begin{enumerate}[label=(\roman*)]
  \item \textbf{Ergebnisstruktur}: DFS erzeugt stets einen \emph{Wald}
        (Menge von Bäumen mit maximaler Tiefe), BFS erzeugt einen oder
        mehrere \emph{flache Bäume} (Tiefe $\leq$ Durchmesser der
        Zusammenhangskomponente, typisch minimal).
  \item \textbf{Komplexität der Ergebnisstruktur}: Der DFS-Wald über
        einem Graphen $G = (V, E)$ hat Tiefe $\Theta(|V|)$ im
        schlechtesten Fall (Pfadgraph), der BFS-Baum hat Tiefe
        $O(\log |V|)$ bei ausgeglichenen Graphen -- eine fundamentale
        Größenordnungsdifferenz.
  \item \textbf{Kantenklassifikation}: DFS unterscheidet vier Kantentypen
        (Baumkanten, Rückwärtskanten, Vorwärtskanten, Querkanten), was eine
        reichere Strukturanalyse erlaubt \cite{Cormen2009}. BFS unterscheidet
        nur Baumkanten und Querkanten.
  \item \textbf{Informationsgehalt}: Die DFS-Stempel (Entdeckungs- und
        Abschlusszeiten) enkodieren die vollständige topologische Struktur
        des Graphen. BFS-Distanzen enkodieren nur metrische, nicht aber
        hierarchische Information.
  \item \textbf{Waldeigenschaft}: DFS erzeugt für beliebige Graphen
        automatisch einen Wald -- dies ist eine inhärente Eigenschaft der
        Rekursionsstruktur. BFS besitzt diese Eigenschaft nicht.
\end{enumerate}
\end{theorem}

\begin{figure}[H]
\centering
\begin{tikzpicture}[
  node distance=1.1cm and 1.5cm,
  box/.style={circle, draw, minimum size=0.75cm, font=\small\ttfamily,
              fill=blue!10},
  hub/.style={circle, draw, minimum size=0.85cm, font=\small\ttfamily\bfseries,
              fill=orange!25},
  treearr/.style={-{Stealth[length=6pt]}, thick, blue!70},
  starr/.style={-{Stealth[length=6pt]}, thick, orange!70!black},
  dashed/.style={-{Stealth[length=6pt]}, dashed, gray}
]

%% ---- DFS-Wald (links) ----
\begin{scope}[xshift=0cm]
  \node[box] (A) at (0, 0)    {$A$};
  \node[box] (B) at (0,-1.4)  {$B$};
  \node[box] (C) at (0,-2.8)  {$C$};
  \node[box] (D) at (0,-4.2)  {$D$};
  \node[box] (E) at (0,-5.6)  {$E$};

  \draw[treearr] (A) -- (B);
  \draw[treearr] (B) -- (C);
  \draw[treearr] (C) -- (D);
  \draw[treearr] (D) -- (E);

  \node[above=0.3cm of A, align=center, font=\small\bfseries]
    {DFS-Wald\\(Tiefenpfad)};
  \node[below=0.2cm of E, font=\scriptsize\color{blue!60}]
    {Tiefe $= 4$, $\Delta = 1$};
\end{scope}

%% ---- BFS-Baum / Sternschema (rechts) ----
\begin{scope}[xshift=6.5cm]
  \node[hub] (H) at (0, -2.8) {$H$};
  \foreach \i/\angle in {1/90, 2/30, 3/-30, 4/-90, 5/150, 6/210}{
    \node[box] (N\i) at ([shift={(\angle:2.4cm)}]H) {$S_{\i}$};
    \draw[starr] (H) -- (N\i);
  }
  \node[above=2.5cm of H, align=center, font=\small\bfseries]
    {BFS-Baum\\(Sternschema)};
  \node[below=2.7cm of H, font=\scriptsize\color{orange!70!black}]
    {Tiefe $= 1$, $\Delta = 6$};
\end{scope}

\end{tikzpicture}
\caption{Gegenüberstellung der durch DFS induzierten Tiefenpfadstruktur
(links, Tiefe $= 4$, $\Delta = 1$) und des durch BFS induzierten
Sternschemas (rechts, Tiefe $= 1$, $\Delta = 6$).
Die Asymmetrie ist offensichtlich: DFS erzeugt einen Wald mit maximaler
Tiefe und minimaler Breite; BFS erzeugt eine flache, sternförmige Struktur.}
\label{fig:dfs_bfs}
\end{figure}

\subsection{Datenbanktheoretische Konsequenz der Asymmetrie}

\begin{corollary}[Datenbanktheoretische Konsequenz]
\label{cor:konsequenz_asymmetrie}
Die algorithmische Asymmetrie von DFS und BFS hat eine direkte
datenbanktheoretische Entsprechung:

\begin{enumerate}[label=(\roman*)]
  \item Der \emph{DFS-Wald} über dem Referenzgraphen eines Schemas ist
        das strukturelle Abbild des \emph{normalisierten Tiefenpfadschemas}.
        Er existiert für beliebige Graphen automatisch und ist azyklisch,
        hierarchisch und eindeutig (bis auf Reihenfolge).
  \item Der \emph{BFS-Baum} über dem Referenzgraphen eines Schemas entspricht
        dem \emph{sternförmigen, denormalisierten Schema}. Er existiert nur
        für Single-Source-Traversierungen und kodiert keine hierarchische
        Tiefenstruktur.
  \item Da DFS automatisch einen \emph{Wald} erzeugt, BFS hingegen
        \emph{keinen Wald}, ist die Ãœberlegenheit des Tiefenpfadschemas
        nicht nur normalisierungstheoretisch, sondern auch algorithmisch
        fundamental: Der DFS-Wald ist die \emph{kanonische, strukturell
        reichere} Darstellung jedes Referenzgraphen.
  \item Die Komplexitätsdifferenz ist qualitativ, nicht nur quantitativ:
        DFS liefert topologische Sortierbarkeit, Zykeldetektierbarkeit
        und vollständige hierarchische Information; BFS liefert keine
        dieser Eigenschaften.
\end{enumerate}
\end{corollary}

\begin{remark}[Praktische Implikation]
Wenn ein Datenbankarchitekt ein Schema entwirft, vollzieht er -- bewusst
oder unbewusst -- eine Traversierung des semantischen Abhängigkeitsgraphen
der Domäne. Wählt er implizit eine DFS-Strategie (erst in die Tiefe einer
Entitätshierarchie, dann zur nächsten), entsteht ein tiefes, normalisiertes
Schema. Wählt er eine BFS-Strategie (alle direkten Attribute einer zentralen
Entität sammeln, ohne in Hierarchien abzusteigen), entsteht ein flaches,
fehleranfälliges Sternschema. Die Asymmetrie der Algorithmen \emph{spiegelt
sich direkt in der Qualität des resultierenden Schemas wider}.
\end{remark}

% ==============================================================
\section{Durchgängiges Praxisbeispiel}
\label{sec:beispiel}
% ==============================================================

\subsection{Domäne: Auftragsabwicklung}

Wir betrachten eine klassische Auftragsabwicklungs-Domäne mit folgenden
Entitäten und natürlichen transitiven Abhängigkeiten:

\[
  \texttt{Bestellung} \xrightarrow{\text{FK}} \texttt{Kunde}
  \xrightarrow{\text{FK}} \texttt{Ort}
  \xrightarrow{\text{FK}} \texttt{Region}
  \xrightarrow{\text{FK}} \texttt{Land}
\]

\subsection{Sternschema (anti-pattern)}

\begin{lstlisting}[caption={Anti-Pattern: Sternschema mit denormalisierten Attributen in Hub-Relation},label=lst:stern]
-- Hub-Relation: enthaelt direkt alle abhaengigen Attribute
CREATE TABLE Bestellung (
    BestellNr     INT PRIMARY KEY,
    Datum         DATE,
    Betrag        DECIMAL(10,2),
    -- Kundenattribute (transitive Abhaengigkeit Stufe 1)
    KundeNr       INT,
    KundeName     VARCHAR(100),
    KundeEmail    VARCHAR(150),
    -- Ortattribute (transitive Abhaengigkeit Stufe 2)
    OrtID         INT,
    OrtName       VARCHAR(80),
    PLZ           CHAR(5),
    -- Regionattribute (transitive Abhaengigkeit Stufe 3)
    RegionID      INT,
    RegionName    VARCHAR(80),
    -- Landattribute (transitive Abhaengigkeit Stufe 4)
    LandID        INT,
    LandName      VARCHAR(60),
    ISO_Code      CHAR(2)
    -- Verletzt 2NF, 3NF und BCNF!
);
\end{lstlisting}

Probleme: Aendert sich der Name eines Kunden, muessen alle seine
Bestellungstupel aktualisiert werden. Existiert ein Kunde noch ohne
Bestellung, kann er nicht eingetragen werden (Einfuegeanomalie).

\subsection{Tiefenpfadschema (BCNF-konform)}

\begin{lstlisting}[caption={Korrekte Tiefenpfadstruktur: BCNF-konformes Schema},label=lst:tief]
CREATE TABLE Land (
    LandID    INT PRIMARY KEY,
    LandName  VARCHAR(60) NOT NULL,
    ISO_Code  CHAR(2)     NOT NULL UNIQUE
);

CREATE TABLE Region (
    RegionID   INT PRIMARY KEY,
    RegionName VARCHAR(80) NOT NULL,
    LandID     INT NOT NULL,
    FOREIGN KEY (LandID) REFERENCES Land(LandID)
);

CREATE TABLE Ort (
    OrtID    INT PRIMARY KEY,
    OrtName  VARCHAR(80) NOT NULL,
    PLZ      CHAR(5),
    RegionID INT NOT NULL,
    FOREIGN KEY (RegionID) REFERENCES Region(RegionID)
);

CREATE TABLE Kunde (
    KundeNr   INT PRIMARY KEY,
    KundeName VARCHAR(100) NOT NULL,
    Email     VARCHAR(150) UNIQUE,
    OrtID     INT NOT NULL,
    FOREIGN KEY (OrtID) REFERENCES Ort(OrtID)
);

CREATE TABLE Bestellung (
    BestellNr INT PRIMARY KEY,
    Datum     DATE         NOT NULL,
    Betrag    DECIMAL(10,2) NOT NULL,
    KundeNr   INT          NOT NULL,
    FOREIGN KEY (KundeNr) REFERENCES Kunde(KundeNr)
);
-- Pfadtiefe: 4 (Bestellung->Kunde->Ort->Region->Land)
-- Jede Relation in BCNF. Keine Anomalien moeglich.
\end{lstlisting}

Der Referenzgraph dieses Schemas ist ein reiner Pfad der Länge 4 --
exakt die durch den DFS-Wald induzierte Tiefenstruktur gemäß
Satz~\ref{thm:dfswald}.

\subsection{Vergleichstabelle}

\begin{table}[H]
\centering
\renewcommand{\arraystretch}{1.35}
\begin{tabular}{lcc}
\toprule
\textbf{Kriterium} & \textbf{Sternschema} & \textbf{Tiefenpfad (BCNF)} \\
\midrule
Normalformgrad            & $<$ 3NF   & BCNF / 4NF \\
Redundanzfaktor $\rho$    & $\gg 1$   & $= 1$ \\
Anomaliefreiheit          & nein      & ja \\
Schreibaufwand (Update)   & $O(n_H)$  & $O(1)$ \\
Suchkosteneffizienz       & $O(n_H)$  & $O(n_{R_d})$ \\
Schemaevolution           & global    & lokal \\
Berechtigungsgranularität & grob      & fein \\
DFS-Wald-Isomorphismus    & nein      & ja \\
Topolog.\ Sortierbarkeit  & begrenzt  & vollständig \\
Zyklusfreiheit            & ungeprüft & inhärent \\
\bottomrule
\end{tabular}
\caption{Zusammenfassender Vergleich Sternschema vs.\ Tiefenpfadschema
nach den formalen Kriterien dieser Arbeit.}
\label{tab:vergleich}
\end{table}

% ==============================================================
\section{Fazit}
\label{sec:fazit}
% ==============================================================

Diese Arbeit hat formal nachgewiesen, dass der tiefenorientierte
Datenbankentwurf dem sternförmigen Entwurf in allen wesentlichen
Gütekriterien überlegen ist. Die Argumentation verlief auf drei
komplementären Ebenen:

\textbf{Normalisierungstheoretisch} (Abschnitte~\ref{sec:normalformen}
und~\ref{sec:anomalien}): Transitive Abhängigkeiten erzwingen tiefe
Pfade, da ihre korrekte Auflösung gemäß Satz~\ref{thm:transitivitaet}
genau eine lineare Zerlegungskette liefert. Sternförmige Schemata
verletzen zwingend die 3NF, sobald transitive Abhängigkeiten in der
Domäne existieren.

\textbf{Algorithmisch} (Abschnitt~\ref{sec:dfs_bfs}): Die Tiefensuche
(DFS) und die Breitensuche (BFS) sind \emph{fundamental asymmetrisch}.
DFS erzeugt stets einen Wald -- eine hierarchische, azyklische,
topologisch sortierbare Struktur mit maximaler Tiefe. BFS erzeugt
keinen Wald, sondern flache, sternförmige Bäume ohne hierarchische
Tiefenstruktur. Diese algorithmische Asymmetrie spiegelt sich direkt
in der Schemaqualität wider: Der DFS-Wald ist die kanonische Darstellung
eines korrekt normalisierten Schemas, der BFS-Baum die Darstellung
eines denormalisierten Sternschemas. Die Folge ist, dass Tiefe \emph{nicht}
nur besser, sondern strukturell kategorial verschieden und
in ihrer Informationsdichte deutlich reicher ist.

\textbf{Praktisch} (Abschnitte~\ref{sec:abfragen},~\ref{sec:wartbarkeit}
und~\ref{sec:beispiel}): Tiefenpfadschemata erzielen überlegene
Selektivität, geringeren Schreibaufwand, feingranulare Zugriffsrechte
und bessere Wartbarkeit.

Das zentrale Prinzip, das aus allen drei Ebenen gleichlautend folgt, lautet:

\begin{quote}
\large\bfseries
Beim Entwurf relationaler Datenbankschemata ist stets die Tiefe der
Breite vorzuziehen. Relationen müssen in die Tiefe gehen, nicht in
die Breite.
\end{quote}

% ==============================================================
% Literatur
% ==============================================================
\newpage
\bibliographystyle{plainnat}
\begin{thebibliography}{99}

\bibitem{Armstrong1974}
Armstrong, W.\,W. (1974).
\newblock Dependency structures of data base relationships.
\newblock \emph{Proceedings of IFIP Congress}, 74, 580--583.

\bibitem{Ambler2003}
Ambler, S.\,W. (2003).
\newblock \emph{Agile Database Techniques: Effective Strategies for the Agile
Software Developer}.
\newblock Wiley.

\bibitem{Boyce1974}
Boyce, R.\,F.; Codd, E.\,F. (1974).
\newblock Relational completeness of data base sublanguages.
\newblock In: Rustin, R. (Hrsg.): \emph{Data Base Systems}, Courant Computer
Science Symposia, Vol.\,6, Prentice-Hall, S.\,65--98.

\bibitem{Codd1970}
Codd, E.\,F. (1970).
\newblock A relational model of data for large shared data banks.
\newblock \emph{Communications of the ACM}, 13(6), 377--387.

\bibitem{Codd1974}
Codd, E.\,F. (1974).
\newblock Recent investigations in relational data base systems.
\newblock \emph{Proceedings of IFIP Congress}, 1021--1036.

\bibitem{Cormen2009}
Cormen, T.\,H.; Leiserson, C.\,E.; Rivest, R.\,L.; Stein, C. (2009).
\newblock \emph{Introduction to Algorithms} (3rd ed.).
\newblock MIT Press.

\bibitem{Date2003}
Date, C.\,J. (2003).
\newblock \emph{An Introduction to Database Systems} (8th ed.).
\newblock Addison-Wesley.

\bibitem{Elmasri2015}
Elmasri, R.; Navathe, S.\,B. (2015).
\newblock \emph{Fundamentals of Database Systems} (7th ed.).
\newblock Pearson.

\bibitem{Fagin1977}
Fagin, R. (1977).
\newblock Multivalued dependencies and a new normal form for relational
databases.
\newblock \emph{ACM Transactions on Database Systems}, 2(3), 262--278.

\bibitem{Graefe1993}
Graefe, G. (1993).
\newblock Query evaluation techniques for large databases.
\newblock \emph{ACM Computing Surveys}, 25(2), 73--169.

\bibitem{Kimball1996}
Kimball, R. (1996).
\newblock \emph{The Data Warehouse Toolkit}.
\newblock Wiley.

\bibitem{Ramakrishnan2003}
Ramakrishnan, R.; Gehrke, J. (2003).
\newblock \emph{Database Management Systems} (3rd ed.).
\newblock McGraw-Hill.

\bibitem{Roddick1995}
Roddick, J.\,F. (1995).
\newblock A survey of schema versioning issues for database systems.
\newblock \emph{Information and Software Technology}, 37(7), 383--393.

\bibitem{Saltzer1975}
Saltzer, J.\,H.; Schroeder, M.\,D. (1975).
\newblock The protection of information in computer systems.
\newblock \emph{Proceedings of the IEEE}, 63(9), 1278--1308.

\bibitem{Tarjan1972}
Tarjan, R.\,E. (1972).
\newblock Depth-first search and linear graph algorithms.
\newblock \emph{SIAM Journal on Computing}, 1(2), 146--160.

\bibitem{Ullman1988}
Ullman, J.\,D. (1988).
\newblock \emph{Principles of Database and Knowledge-Base Systems}, Vol.\,1.
\newblock Computer Science Press.

\end{thebibliography}

\end{document}
